\section{A local winding count}\label{sec:localization}

We now combine the geometry of Section~\ref{sec:sectors} with the
diagram statements of Section~\ref{sec:diagrams}.
Continue to assume the action and strict inequalities of
Proposition~\ref{prop:triangle-obstruction}, and retain the sphere maps
\(f_n\) and the constant \(\epsilon>0\) from
Lemma~\ref{lem:outer-cap}.

Write $\Sigma$ for the parameter sphere; for the geometric map
$f_n:\Sigma\to Y$, Lemma~\ref{lem:controlled-diagrams} supplies a
separate labelled diagram $\Sigma\to\mathcal K$ on that sphere.
Projected mirror coordinates of its group-labelled vertices, extended
affinely over parameter triangles, will define a planar map
$q_n:\Sigma\to\mathbb R^2$ that approximates coordinates on a small
patch and has nonzero local winding.  Separation in $Y$ will force every
occurrence of a contributing target triangle into the patch, making that
winding a weighted sum of whole-sphere signed counts, which
Proposition~\ref{prop:sphere-counts} makes zero.

\subsection{A patch and its coordinates}\label{sec:local-patch}

Choose one of the six sectors, say the sector for
\(U=u^2,H=z_{uv}\).  Its exponent coordinates form the nonnegative
quadrant in a Euclidean plane.  Choose a small closed coordinate
rectangle
\[
 B=[\lambda_0,\lambda_1]\times[\mu_0,\mu_1]
\]
in the interior of this sector, with \(\lambda_0,\mu_0>0\), and so small
and close to the origin that all its radii satisfy
\begin{equation}\label{eq:small-patch}
 0<r(p)<\min(\epsilon/3,1/4)\qquad(p\in B).
\end{equation}
Write \(\kappa(p)=(\lambda,\mu)\) on \(B\).  Fix an interior point
\(p_*\) and a smaller closed coordinate rectangle
\(B'\subset\operatorname{int}B\) containing \(p_*\) in its interior.

We regard \(P\) and the second disk as a fixed triangulated parameter
sphere, with \(B\) a subcomplex.  One explicit way to arrange this is
as follows.  In each sector take the triangle whose two outer vertices
lie on its boundary rays at radius one.  The six triangles form a
polygonal cone disk \(Q\).  On a radial line, its outer radius \(\rho\)
is at least \(\cos(\theta/2)>\sqrt3/2\).  Map \(Q\) to \(P\) by a
radial homeomorphism that fixes radii at most \(1/2\) and is linear
from radius \(1/2\) to the outer radius, sending \(\rho\) to \(1\).
These maps agree on the common rays.  The rectangle \(B\) is fixed
pointwise.  Extend its four straight sides within its triangle to
partition that triangle into finitely many convex polygons; triangulate
each polygon and the other five triangles.  Finally cone the resulting
boundary subdivision to a new parameter vertex for the second disk.
This supplies the required finite parameter triangulation.

For \(\xi=\sum_D\xi_D e_D\), define the bounded linear map
\begin{equation}\label{eq:coordinate-projection}
 J\xi=\left(\frac{\xi_{D_u}-\xi_{D_{uv}}}{2},
                 \frac{\xi_{D_{uv}}}{2}\right).
\end{equation}
By \eqref{eq:mirror-translations},
\begin{equation}\label{eq:projection-identity}
 Jv_*(p)=\kappa(p)
\end{equation}
throughout the chosen sector.  This is a coordinate projection on that
sector; no global injectivity of \(J\) is asserted.

Apply Lemma~\ref{lem:controlled-diagrams} to each sphere map \(f_n\),
preserving \(B\) as a subcomplex.  We obtain a subdivided combinatorial
sphere map to \(\mathcal K\).  Denote the group label at a vertex by
\(s\).  There is one constant \(C_1\), independent of \(n\), for which
\begin{equation}\label{eq:local-all-points}
 d_Y(sy,f_n(p))\le C_1
\end{equation}
whenever \(s\) is a vertex label of a final parameter triangle and
\(p\) is any point of that triangle.

On this parameter sphere, assign to the vertex with label \(s\) the
plane value \(JV(s)/n\), and extend affinely on each abstract triangle.
Call the resulting continuous map \(q_n\).  Affineness refers to the
abstract triangle coordinates, regardless of the homeomorphisms used
to place the triangles in the parameter sphere.

\begin{lemma}\label{lem:projected-approximation}
On \(B\), the maps \(q_n\) converge uniformly to \(\kappa\).
\end{lemma}

\begin{proof}
Fix \(p\in B\) and any vertex label \(s\) of a triangle containing \(p\).
The bounds \eqref{eq:local-all-points} and \eqref{eq:grid-error} imply
\[
 d_Y(sy,t_n(p)y)\le C_1+C_0.
\]
By \eqref{eq:orbit-control} and \eqref{eq:mirror-lipschitz}, there is a
constant \(C_2\), independent of \(n,p,s\), such that
\[
 \|V(s)-V(t_n(p))\|_1\le C_2.
\]
The second convergence in \eqref{eq:radial-coordinates} has error
\(O(1/n)\), so
\[
 \left\|\frac{V(s)}n-v_*(p)\right\|_1=O(1/n)
\]
uniformly for every such vertex label.  Apply \(J\) and
\eqref{eq:projection-identity}.  Since \(q_n(p)\) is a convex combination
of those vertex values, it satisfies the same uniform error bound from
\(\kappa(p)\).  This is precisely where the all-points form of
\eqref{eq:local-all-points} is used.
\end{proof}

\subsection{The winding calculation}

For completeness, we spell out the elementary planar count we use.
For a loop avoiding a point \(w\), choose continuous local arguments
of its vectors from \(w\) on successive short parameter intervals.
Matching the arguments at their common endpoints gives a total angle
change in \(2\pi\mathbb Z\).  Its quotient by \(2\pi\) is the winding
number.  It is unchanged under a homotopy avoiding \(w\): local angle
choices vary continuously, and a continuous integer-valued function is
constant.

If a finite oriented triangulated disk is mapped affinely to the plane,
and \(w\) avoids all projected edges, sum the winding numbers of its
oriented triangle boundaries.  Interior edges occur twice with opposite
orientations and their angle increments cancel.  The sum is therefore
the winding number of the outer boundary.  For one affine triangle the
term is zero if \(w\) is outside its convex hull, and is its orientation
sign if \(w\) lies in its interior.  In the first case one may contract
the loop inside its convex hull; in the second, rays from \(w\) meet
the triangle boundary once and its direction makes one full turn.
A degenerate triangle contributes zero because \(w\) avoids its
projected edges.

Choose a small open plane ball \(W\) about \(\kappa(p_*)\) whose closure
lies in \(\kappa(\operatorname{int}B')\).  By
Lemma~\ref{lem:projected-approximation}, for large \(n\) and every \(w\)
in a sufficiently small fixed such ball, the following hold:
\begin{enumerate}
\item the straight-line homotopy from \(q_n|_{\partial B}\) to
      \(\kappa|_{\partial B}\) avoids \(w\);
\item every \(p\in B\) with \(q_n(p)=w\) lies in
      \(\operatorname{int}B'\).
\end{enumerate}
Indeed, take the approximation error and the radius of \(W\) smaller
than half the positive distance from \(\kappa(p_*)\) to
\(\kappa(B\setminus\operatorname{int}B')\).
For each large \(n\), choose \(w\in W\) avoiding every projected edge
of the entire finite sphere diagram.  A finite union of line segments
cannot contain an open ball, so this choice is possible.

The positively or negatively oriented rectangle \(\kappa(\partial B)\)
has winding number \(1\) or \(-1\) about \(w\).
The preceding homotopy and cancellation argument show that the sum of
signed coverings of \(w\) by the triangles of \(B\) is nonzero.
Every contributing triangle is noncollapsed, has a nondegenerate planar
image, and contains a parameter point of \(B'\) projecting to \(w\).

\subsection{All occurrences of a contributing cell are local}

\begin{lemma}\label{lem:localized-degree}
For all sufficiently large \(n\), every occurrence in the sphere
diagram of a target triangle contributing to the winding count on
\(B\) lies entirely in \(\operatorname{int}B\).
\end{lemma}

\begin{proof}
By \eqref{eq:region-separation}, the \(f_n\)-images of \(B'\) and
\(P\setminus\operatorname{int}B\) are separated by a positive multiple
of \(n\).  On the second disk,
Lemma~\ref{lem:outer-cap} gives distance at least \(\epsilon n\) from
\(y\).  On \(B'\), \eqref{eq:small-patch} and
\eqref{eq:radial-coordinates} give distance at most
\((\epsilon/3)n+o(n)\) from \(y\).  Thus the images of \(B'\) and the
entire complement of \(\operatorname{int}B\) in the sphere have mutual
distance tending to infinity, uniformly.

Let \(\tau\) be a contributing labelled target triangle.  One of its
occurrences contains \(p\in B'\) with \(q_n(p)=w\).
Fix any one of the three group vertex labels \(s\) of \(\tau\).
Every other occurrence of this same target cell has the same vertex
label \(s\), including the full translate information.
If some occurrence met the complement of \(\operatorname{int}B\),
take a point \(p'\) in that intersection.  The all-points bound
\eqref{eq:local-all-points} would give
\[
 d_Y(f_n(p),f_n(p'))\le
 d_Y(f_n(p),sy)+d_Y(sy,f_n(p'))\le2C_1,
\]
contradicting the uniform separation for large \(n\).
This also excludes an occurrence touching \(\partial B\).
\end{proof}

\begin{proof}[Proof of Proposition~\ref{prop:triangle-obstruction}]
Choose \(n\) large enough for the constructions and
Lemma~\ref{lem:localized-degree}.  Orient every labelled target triangle
of \(\mathcal K\) once and for all.  Its three group vertices have fixed
projected values \(JV(s)/n\), so its affine planar image is fixed.
If it covers \(w\), its determinant has a fixed nonzero sign
\(\eta_\tau\in\{1,-1\}\).  The covering sign of each occurrence is
\(\eta_\tau\) times that occurrence's orientation sign relative to the
chosen orientation of the target.

Group the nonzero winding sum on \(B\) by its contributing target
triangles.  By Lemma~\ref{lem:localized-degree}, every occurrence of
each of these target triangles in the whole sphere lies in \(B\).
Its total contribution is therefore
\[
 \eta_\tau\,
 \bigl(\text{signed number of occurrences of }\tau
       \text{ in the sphere}\bigr)=0
\]
by Proposition~\ref{prop:sphere-counts}.  Summing gives zero, contrary
to the winding number \(1\) or \(-1\).
Distinct target cells remain separate in this sum even if their
projected vertex data agree.  The contradiction proves the proposition.
\end{proof}
